\section{Introduction}

Plagiarism undermines academic integrity and poses significant challenges to educational institutions globally. The prevalence of digital content and sophisticated paraphrasing techniques have rendered traditional plagiarism detection methods—based on string matching and hash-based fingerprinting—increasingly ineffective.

Consider the evolution of plagiarism detection: Early systems relied on hash-based matching, which effectively detects near-duplicate content but fails against lexical variation and semantic paraphrasing. Recent advances in machine learning and deep learning offer promising alternatives, yet comprehensive comparative evaluation remains limited.

\subsection{Motivation}

This research is motivated by three key observations:

\begin{enumerate}
  \item \textbf{Technology gap} --- Advances in transformer-based models (BERT) promise superior detection, yet cost-benefit analysis remains unexplored
  
  \item \textbf{Institutional constraints} --- Educational institutions face diverse resource constraints, requiring multiple detection approaches rather than one-size-fits-all solutions
  
  \item \textbf{Evaluation gap} --- Existing literature emphasizes individual approaches without systematic comparative evaluation providing practical deployment guidance
\end{enumerate}

\subsection{Research Contributions}

This paper provides:

\begin{enumerate}
  \item \textbf{Systematic comparative evaluation} of three detection paradigms: TF-IDF (classical ML), BiLSTM (deep learning), and BERT (transformers)
  
  \item \textbf{Quantitative analysis} of accuracy-efficiency trade-offs, decomposed by plagiarism type and deployment scenario
  
  \item \textbf{Practical deployment recommendations} aligned with institutional constraints and performance requirements
  
  \item \textbf{Comprehensive evaluation framework} for plagiarism detection reproducible across implementations
\end{enumerate}

\subsection{Paper Organization}

Remaining sections are organized as follows:
\begin{itemize}
  \item Section \ref{sec:related_work} reviews existing plagiarism detection approaches
  \item Section \ref{sec:problem} articulates specific research objectives
  \item Section \ref{sec:dataset} describes dataset construction and characteristics
  \item Section \ref{sec:methodology} presents three detection approaches in detail
  \item Section \ref{sec:implementation} details software implementation
  \item Section \ref{sec:experimental} describes experimental design and methodology
  \item Section \ref{sec:results} presents experimental results
  \item Section \ref{sec:discussion} analyzes findings and implications
  \item Section \ref{sec:conclusion} summarizes contributions and limitations
  \item Section \ref{sec:future} outlines promising future research directions
\end{itemize}
